% Generated draft: requires mathematical and source review.
\documentclass[11pt]{article}
\usepackage[margin=25mm]{geometry}
\usepackage{fontspec}
\setmainfont{Latin Modern Roman}
\usepackage{amsmath,amssymb,mathtools,unicode-math}
\setmathfont{Latin Modern Math}
\usepackage{enumitem,needspace,xurl,hyperref}
\setlength{\parindent}{0pt}
\setlength{\parskip}{5pt}
\setlength{\emergencystretch}{3em}
\newcommand{\N}{\mathbb{N}}
\newcommand{\Z}{\mathbb{Z}}
\newcommand{\Q}{\mathbb{Q}}
\newcommand{\R}{\mathbb{R}}
\newcommand{\C}{\mathbb{C}}
\newcommand{\F}{\mathbb{F}}
\newcommand{\A}{\mathbb{A}}
\newcommand{\PP}{\mathbb{P}}
\newcommand{\E}{\mathbb{E}}
\newcommand{\eps}{\varepsilon}
\newcommand{\dd}{\,\mathrm{d}}
\begin{document}
\section*{UnsolvedMath ID: 2500004}
\textbf{Source problem code:} catalogue entry
\subsection{Definitions and mathematical statement}
Let $B$ be a fixed obstacle in $\R^3$ and let $\Omega=\R^3\setminus B$ be the exterior fluid region. The source specifies viscous incompressible flow with
\[
 u(x,t)\longrightarrow U_\infty\ne0\qquad(|x|\to\infty).
\]
A customary model, per unit fluid density, is
\[
 \partial_tu+(u\cdot\nabla)u=-\nabla p+\nu\Delta u,
 \qquad\nabla\cdot u=0\quad\text{in }\Omega,
\]
with viscosity $\nu>0$. The source does not supply obstacle regularity, a boundary condition on $\partial B$, a solution class, or whether the problem is steady or time dependent. In particular, no-slip $u=0$ on the fixed obstacle is a common convention but is not imposed here without source support.

With $D(u)=\tfrac12(\nabla u+\nabla u^{\mathsf T})$, the standard total viscous dissipation rate per unit density is
\[
 \varepsilon(t)=2\nu\int_{\Omega}|D(u)(x,t)|^2\,dx,
\]
when finite. A stationary rate or long-time averaged rate
\[
 \overline\varepsilon=\limsup_{T\to\infty}\frac1T\int_0^T\varepsilon(t)\,dt
\]
are distinct possible targets; the source does not choose between them or an instantaneous bound with initial-data dependence.

\textbf{Problem.} Obtain realistic rigorous upper bounds on the energy-dissipation rate for this three-dimensional exterior flow past the fixed obstacle with nonzero uniform velocity at infinity. Once the setting is fixed, a typical goal is an explicit bound in terms of viscosity, incoming speed, and obstacle geometry. If a characteristic length $L$ is chosen, one can express sharpness through $\operatorname{Re}=|U_\infty|L/\nu$ and a dimensionless rate $\overline\varepsilon/(|U_\infty|^3L^2)$, but the source prescribes neither $L$ nor a Reynolds-number scaling. ``Realistic'' therefore needs a specified physical or asymptotic criterion, not an invented numerical bound.

The supplied research summary reports rigorous drag/dissipation bounds in special geometries and says fully general physically sharp bounds for arbitrary obstacles remain partially open. It gives no precise results or citations. This is retained as an attributed progress assessment rather than a verified theorem or current literature conclusion.
\subsection{Short English statement}
Find physically sharp upper bounds for dissipation in viscous incompressible flow past a fixed three-dimensional obstacle with nonzero uniform velocity at infinity. The boundary condition, solution class, averaging convention, geometry class, and meaning of realistic sharpness must be specified before a unique bound can be formulated.
\subsection{Sources}
\begin{itemize}
\item Local UnsolvedMath registry, permanent ID 2500004, source code catalogue entry, ``Dissipation bounds for flow past an obstacle''; supplied status partially solved, classification PARTIAL-PROGRESS, updated 2026-08-06.
\item Constantin, \emph{Problems} (2001), source-order Question 4, \url{https://web.math.princeton.edu/~const/2k.pdf}. The record reports PDF-text extraction and access on 2026-07-29. It says the source presents an open item and current status requires release review. Rights: NEEDS\_REVIEW; public author-hosted PDF.
\item The supplied summary mentions special-geometry dissipation/drag estimates but supplies no further references or theorem statements.
\end{itemize}
\end{document}
